% ============================================================================
%  Actividad 5 · Anexo B. Manual de uso del prototipo
%
%  Propiedad de US-AVANCE-5, sub-tarea F. Un solo archivo que se compone dos
%  veces: como Anexo B de la Parte V de main_completo.tex y como PDF propio
%  desde main_manual.tex. Un parrafo vive una vez.
%
%  Fuentes de verdad del contenido, ninguna reescrita aqui:
%    frontend/app/utils/navegacion.ts        rutas, modulos y facetas
%    frontend/app/utils/permisos.generated.ts  rol minimo por ruta y por rama
%    docs/security.md                        matriz de permisos y jerarquia
%    docs/espacios-de-trabajo.md             composicion de inicio por rol
%    frontend/i18n/locales/es.json           rotulos reales de la interfaz
%    contenido/a4_02_prototipos.tex          vocabulario de las siete pantallas
%    contenido/a5_04_usabilidad.tex          tareas T1 a T6 del 20-ago-2026
%
%  Las capturas de las siete pantallas salen de figuras/a4/tema/ porque son el
%  estado entregado del portal, el mismo chasis que muestra la captura del
%  tablero tomada contra el despliegue. Las dos figuras de figuras/a5/ se
%  capturaron contra Cloud Run con sesion de perfil analista.
%
%  Cada pantalla abre pagina: la captura ocupa cerca de media caja de texto y
%  sin el salto el flotante aterriza junto al titular de la pantalla siguiente.
% ============================================================================

\uxsection{Anexo B. Manual de uso del prototipo}
\phantomsection
\label{sec:a5-manual}

Este anexo explica cómo se usa el prototipo de Karisma Data. Está escrito para quien lo abre por
primera vez. Describe los cuatro perfiles, recorre las siete pantallas en el orden en que se
navegan y cierra con lo que aparece cuando algo no sale. Cada pantalla termina con un ejemplo
completo, y ninguno de esos ejemplos es inventado: los siete salen de las tareas que cinco
personas ejecutaron en la prueba de usabilidad del jueves 20 de agosto de 2026.

% ============================================================================
\subsection{Qué es y qué no es este prototipo}

\begin{uxdestacado}
\textbf{Karisma Data es un prototipo de alta fidelidad operable con datos sintéticos. No está
conectado a sistemas reales de ninguna institución.} Ninguna cifra que aparece en pantalla mide
un negocio real y ninguna sostiene una decisión.
\end{uxdestacado}

El portal declara esa condición él mismo, sin que haya que buscarla. Una franja de alcance
acompaña a las diez rutas y repite la misma frase: \enquote{Prototipo de alta fidelidad de
Karisma Data con datos sintéticos. No está conectado a sistemas reales de ninguna institución.}
Los bloques que muestran listas o cifras de muestra llevan además su propio distintivo,
\enquote{Datos de ejemplo} o \enquote{Datos sintéticos}, en su encabezado.

\siHayDemoWeb
  {El portal se abre en un navegador de escritorio, en \urlDemoWeb. No hace falta instalar nada.}
  {El portal se abre en un navegador de escritorio y no hace falta instalar nada:
   \urlDemoWeb.}
La interfaz es bilingüe y se cambia de idioma en cualquier momento con los dos botones
\textbf{ES} y \textbf{EN} de la cabecera; este manual cita el rótulo en español y, cuando ayuda a
localizarlo, el equivalente en inglés entre paréntesis.

\begin{uxlist}
  \lead{Qué sí es}{un producto navegable de principio a fin. La sesión es real, el permiso por
    rol se comprueba en cada petición, la búsqueda del catálogo consulta al servidor, la
    exportación produce un archivo con un enlace firmado y el tablero dibuja una serie
    preagregada de verdad.}
  \lead{Qué no es}{un sistema en producción. Los datos provienen de generadores sintéticos con
    semilla fija, las proyecciones son rectas ajustadas a esa historia sintética y las
    respuestas del asistente están escritas de antemano. La pantalla del asistente y la de
    exportación lo declaran en su propia franja.}
  \lead{Qué falta por construir}{tres de las cuatro ramas de Administración
    ---solicitudes y aprobaciones, bitácora de accesos e integraciones--- se nombran en la
    pantalla y no están construidas. La propia pantalla lo dice: \enquote{Las solicitudes, la
    bitácora de accesos y las integraciones llegan después de esta entrega.}}
\end{uxlist}

La cabecera termina con una insignia \textbf{Entorno} que declara sobre qué copia del portal se
está trabajando. En la dirección pública dice \texttt{demo}; en las capturas tomadas sobre el
entorno de desarrollo dice \texttt{local}.

% ============================================================================
\subsection{Los cuatro perfiles y lo que ve cada uno}

El portal tiene cuatro perfiles y ninguno más: \textbf{operativo}, \textbf{analista},
\textbf{directivo} y \textbf{administrador}. Forman un orden total, de modo que un perfil alcanza
todo lo que alcanza el rango inferior. Por eso la última columna de la tabla dice qué añade cada
uno respecto del anterior, y no qué le está vedado a los demás.

\begin{uxtablalarga}{|L{2.2cm}|L{3.4cm}|L{4.0cm}|Y|}{Los cuatro perfiles: para quién es cada uno, qué módulos abre y qué añade}{\thd{Perfil} & \thd{Para quién es} & \thd{Módulos que abre} & \thd{Qué añade respecto del anterior}}
  Operativo \newline \emph{(Operations)} & Quien localiza una cifra y la comprueba: consulta puntual, control y auditoría & Inicio, Exploración y extracción, Gobierno del dato y Asistente conversacional & Es el rango mínimo del portal. El catálogo y el diccionario abren con cualquier sesión; el asistente pide ya este perfil \\
  Analista \newline \emph{(Analyst)} & Quien construye consultas, cruza fuentes y extrae datos para trabajar con ellos & Los anteriores más Tableros e indicadores y Exportaciones & Las dos continuaciones gobernadas del módulo de exploración: la serie preagregada del tablero y la solicitud de exportaciones \\
  Directivo \newline \emph{(Executive)} & Quien pide un estado, no un conjunto de datos: seguimiento y decisión & Los mismos que el analista & No abre ninguna pantalla nueva. Cambia la composición de Inicio, que arranca con los indicadores del último corte, y es el rango que la matriz reserva para los resúmenes ejecutivos, cuya ruta está declarada y todavía no montada \\
  Administrador \newline \emph{(Administrator)} & Quien administra las cuentas y los permisos del portal & Todos los anteriores más Administración & Consultar las cuentas del portal, cambiar el rol de una cuenta y desactivarla o reactivarla \\
\end{uxtablalarga}

Dos reglas que conviene conocer antes de entrar, porque explican casi todo lo que el portal hace
al cambiar de perfil:

\begin{uxlist}
  \lead{Lo que un perfil no puede abrir se nombra, no se esconde}{la pantalla de exploración
    nombra sus dos continuaciones gobernadas, declara qué perfil piden y las deja visibles sin
    ser accionables. Ocultarlas impediría enterarse de que existen y de que se pueden pedir.}
  \lead{Cada perfil abre un espacio de trabajo distinto, y eso es el diseño}{Inicio se compone de
    tres maneras según quien entra. La operativa abre con el buscador, la del analista con las
    consultas guardadas y las exportaciones, y la directiva con los indicadores del último corte.
    No es una limitación del prototipo: es la decisión de arrancar con tres composiciones
    sensatas en lugar de imponer una vista única a ocho perfiles distintos.}
\end{uxlist}

% ============================================================================
\clearpage
\subsection{Cómo entrar}

La pantalla de acceso ofrece dos caminos y recomienda uno. El recomendado es el bloque
\textbf{Entrar como perfil de demostración} (\emph{Sign in as a demo profile}): cuatro botones, uno
por perfil, que abren sesión sin contraseña y declaran a dónde llevan. El otro es el formulario de
\textbf{usuario y contraseña}, plegado bajo su propio disparador; en este prototipo nadie tiene una
cuenta con contraseña, y la pantalla lo advierte.

\figuraux[0.56\linewidth]{figuras/a4/tema/institucional_claro_0_acceso.png}{Pantalla de acceso. Los
cuatro botones de perfil ---Operativo, Analista, Directivo y Administración--- ocupan la superficie
principal, cada uno con la frase que declara a dónde lleva, y el formulario de usuario y contraseña
queda plegado debajo.}{fig:a5-manual-acceso}

\begin{uxpreg}
  \item Abre la dirección pública del portal. Si no hay sesión, cualquier ruta protegida devuelve
    a esta pantalla.
  \item Pulsa el perfil con el que quieres entrar: Operativo, Analista, Directivo o
    Administración.
  \item El portal abre la pantalla inicial de ese perfil. Los tres primeros llegan a Inicio, con
    su composición; Administración llega directamente a la pantalla de administración.
  \item Para cambiar de perfil no hace falta cerrar sesión: el selector de la cabecera, que
    muestra el perfil en curso, abre los cuatro y cambia de sesión en el sitio.
\end{uxpreg}

El rótulo ámbar \enquote{Acceso de demostración, habilitado únicamente en este prototipo} señala
exactamente eso: la puerta sin contraseña existe aquí y no es un comportamiento que el producto
proponga para un entorno real. La sesión dura 30 minutos; cuando caduca, la pantalla de acceso lo
explica y el formulario de credenciales se despliega solo.

Todas las pantallas comparten el mismo chasis, así que lo aprendido en una sirve en las demás:
barra lateral con los módulos del perfil y, bajo el activo, sus ramas; cabecera con el buscador,
\textbf{Apariencia} ---tema y modo---, perfil, idioma y entorno; y, como primera parada del
tabulador, un enlace para \textbf{saltar al contenido}. El asistente cuelga al pie de la barra
lateral: es transversal a los cuatro módulos y no pertenece a ninguno.

% ============================================================================
\clearpage
\subsection{Inicio}

\textbf{Para qué sirve.} Es el punto de partida de la sesión. Reúne en una sola pantalla el
buscador unificado, las búsquedas recientes, las fuentes favoritas, las alertas y el perfil, y los
ordena según quien entró.

\textbf{Qué se ve.} Un saludo con el nombre y el perfil, una frase que explica por qué el espacio
abre por donde abre y, debajo, los bloques de la composición. Cada bloque con contenido lleva su
distintivo \enquote{Datos de ejemplo}, y las listas usan el vocabulario real del catálogo: los
nombres físicos \texttt{ratio\_lcr}, \texttt{sdo\_cap}, \texttt{dias\_mora},
\texttt{nocional\_usd} y \texttt{mto\_disp}, y las fuentes \texttt{liquidez}, \texttt{creditos},
\texttt{derivados} y \texttt{regulatorio}.

\figuraux[0.92\linewidth]{figuras/a4/tema/institucional_claro_1_inicio.png}{Inicio con sesión de
perfil operativo. El buscador unificado domina la pantalla porque localizar una cifra y comprobarla
es la acción que este perfil siempre ejecuta. Debajo, las búsquedas recientes y las fuentes
favoritas, ambas marcadas como datos de ejemplo, y las alertas con su nivel de severidad. La barra
lateral despliega las cinco ramas del módulo activo.}{fig:a5-manual-inicio}

\begin{uxpreg}
  \item Escribe en \textbf{Buscar en el catálogo} un nombre físico como \texttt{ratio\_lcr} o un
    concepto de negocio como \enquote{cobertura de liquidez}, y pulsa \textbf{Buscar}.
  \item Revisa \textbf{Búsquedas recientes} y \textbf{Fuentes favoritas} para retomar algo ya
    consultado sin volver a teclearlo.
  \item Atiende \textbf{Mis alertas}, que marca cada señal como Alta, Media o Informativa.
  \item Consulta \textbf{Mi perfil} para confirmar con qué rol estás trabajando y cuándo fue tu
    último acceso. Este bloque solo aparece en la composición operativa.
\end{uxpreg}

\begin{uxnota}[Ejemplo, tarea T1 de la prueba del 20 de agosto de 2026]
La tarea pedía \emph{entrar como analista y confirmar que se alcanzó el espacio correspondiente}.
Los dos participantes del bloque estandarizado pulsaron el botón \textbf{Analista} de la pantalla
de acceso y reconocieron su espacio porque abre con \textbf{Consultas guardadas} y
\textbf{Exportaciones recientes} en lugar del buscador. Ambos lo completaron sin errores, en
24 y 29 segundos frente a una meta de 45, y calificaron la facilidad con el máximo de la escala,
7 sobre 7. Inicio no tuvo tarea cronometrada propia: lo que esta pantalla aportó al estudio es el
destino de T1 y el hallazgo H2, según el cual el punto de partida de la búsqueda no siempre resulta
evidente ---tres de los cinco participantes probaron primero el buscador de la cabecera---.
\end{uxnota}

% ============================================================================
\clearpage
\subsection{Exploración y extracción}

\textbf{Para qué sirve.} Es la pantalla de descubrimiento: encontrar el campo antes de pedir el
dato. Responde a la pregunta \enquote{¿qué dato necesito y dónde está?}.

\textbf{Qué se ve.} Tres zonas. Arriba, el buscador \textbf{¿Qué dato necesitas?}; a la izquierda,
la columna \textbf{Dominios} con el conteo de coincidencias de cada uno; a la derecha, la tabla
\textbf{Campos del catálogo} con seis columnas: nombre físico, nombre de negocio, dominio, fuente,
responsable y certificación. La tabla declara cuántos campos muestra de cuántos coinciden, anuncia
por qué columna está ordenada y distingue los tres estados de certificación ---Certificado, En
revisión y Obsoleto--- por color y por icono, nunca por color solo.

\figuraux[0.92\linewidth]{figuras/a4/tema/institucional_claro_2_exploracion.png}{Exploración y
extracción con sesión de perfil analista, abierta en frío desde una dirección que ya lleva el
término \texttt{saldo}. La columna de dominios acota sobre el total de coincidencias y la tabla
densa publica sus seis columnas; cada fila lleva el estado de certificación con su icono, de modo
que Certificado, En revisión y Obsoleto se distinguen sin depender del
color.}{fig:a5-manual-exploracion}

\begin{uxpreg}
  \item Escribe al menos dos caracteres y envía el formulario. La consulta sale al enviar, no con
    cada tecla, y el término queda en la dirección: un enlace compartido abre la pantalla ya
    filtrada.
  \item Acota con un dominio de la columna izquierda. El filtro no vuelve a pedir la página.
  \item Lee en la fila la fuente, el responsable y la certificación del campo.
  \item Pulsa el nombre físico de la fila para abrir el \textbf{linaje} del campo, que se dibuja
    en el mismo panel que usa Gobierno del dato.
  \item Continúa con el resultado hacia la ficha completa en Gobierno, hacia el tablero o hacia la
    exportación. El bloque \textbf{Llevar el resultado a} reúne esas salidas y declara qué perfil
    pide cada una.
\end{uxpreg}

\begin{uxnota}[Ejemplo, tarea T2 de la prueba del 20 de agosto de 2026]
La tarea pedía \emph{encontrar el campo de cobertura de liquidez e identificar su fuente, su
responsable y su certificación}. Los dos participantes escribieron el concepto de negocio en el
buscador de la pantalla, localizaron el campo en la tabla y leyeron los tres metadatos en su
propia fila. Los dos terminaron con éxito, en 51 y 62 segundos frente a una meta de 90, y los dos
cometieron el mismo error: probaron primero el buscador de la cabecera. Es la evidencia del
hallazgo H2 y la razón de la recomendación R2, que pide diferenciar la búsqueda global de la
contextual. La facilidad percibida fue de 6 sobre 7 en ambos casos.
\end{uxnota}

% ============================================================================
\clearpage
\subsection{Tableros e indicadores}

\textbf{Para qué sirve.} Es la zona del módulo de exploración que muestra la evolución de una
métrica en el tiempo y su proyección al mes siguiente. Pide perfil analista.

\textbf{Qué se ve.} Dos bloques, uno encima del otro. Arriba, \textbf{Proyección del próximo mes}
con distintivo \enquote{Datos sintéticos} y tres tarjetas: Cobertura de liquidez (LCR),
Saldo disponible promedio diario y Concentración en divisa extranjera. Cada tarjeta publica su
cifra, su variación contra el mes anterior y el método que la produjo, con su R\textsuperscript{2}
a la vista. Abajo, la \textbf{Serie preagregada de liquidez}, con los controles de Métrica,
Agrupación y Densidad de lectura, la gráfica, su leyenda, un resumen escrito que dice lo mismo que
la gráfica para quien usa lector de pantalla, y el bloque \textbf{Procedencia del dato}.

\figuraux[0.56\linewidth]{figuras/a5/tableros_analista.png}{Tableros e indicadores con sesión de
perfil analista sobre el portal desplegado. Arriba, las tres tarjetas de proyección con su cifra,
su variación frente a junio de 2026 y el método declarado en la propia tarjeta. Abajo, la serie
preagregada con sus controles de métrica, agrupación y densidad, la leyenda de cinco líneas, el
resumen escrito equivalente a la gráfica y la procedencia del dato, que publica el archivo de
origen, la semilla y las transformaciones.}{fig:a5-manual-tableros}

\begin{uxpreg}
  \item Elige la \textbf{Métrica} ---Saldo disponible, Razón LCR o Número de posiciones--- y la
    \textbf{Agrupación} ---Unidad de negocio, Divisa, Bucket de vencimiento o Serie individual---.
  \item Ajusta la \textbf{Densidad de lectura}: Resumen, Detalle o Carga completa. La última dibuja
    los 500\,000 puntos preagregados y existe como evidencia de rendimiento.
  \item Usa la barra inferior para acercar un tramo, o \textbf{Ver todo el periodo} para volver.
  \item Pulsa \textbf{Ver la tabla de detalle} si prefieres los números a la gráfica: toda gráfica
    del portal tiene su alternativa en tabla.
  \item Abre \textbf{Procedencia del dato} para ver el archivo de origen, las filas crudas y las
    preagregadas, la semilla del generador y las transformaciones aplicadas.
  \item Pulsa \textbf{Ver la serie} en una tarjeta de proyección para desplegar su gráfica mensual
    con el tramo proyectado.
\end{uxpreg}

\begin{uxdestacado}
\textbf{Cómo se distingue un dato de una proyección.} La tarjeta de la cobertura de liquidez marca
130.2\,\% con una variación de +0.24\,\% frente a junio de 2026, y declara su método: proyección
lineal por mínimos cuadrados sobre 12 meses de datos sintéticos, horizonte de un mes,
R\textsuperscript{2} de 0.02. Al abrir la serie, la leyenda separa \enquote{Observado, línea
continua} de \enquote{Proyectado, línea discontinua y marcador hueco}, el rótulo
\enquote{Último mes observado: jun 2026} fija dónde termina lo medido, y el origen del dato nombra
el campo \texttt{ratio\_lcr} de la serie preagregada. La frase final no deja lugar a la
ambigüedad: \enquote{No hay aprendizaje automático detrás de esta cifra: es una recta ajustada a
la historia sintética del prototipo.}
\end{uxdestacado}

\figuraux[0.88\linewidth]{figuras/a5/serie_observado_proyectado.png}{Detalle que abre el botón
\enquote{Ver la serie} de la tarjeta de cobertura de liquidez. El último punto se dibuja con línea
discontinua y marcador hueco, la leyenda nombra los dos tramos, el rótulo declara cuál fue el
último mes observado y el bloque de origen publica el campo, la agregación y el método con su
R\textsuperscript{2}.}{fig:a5-manual-serie}

\begin{uxnota}[Ejemplo, tarea T4 de la prueba del 20 de agosto de 2026]
La tarea pedía \emph{explicar la tendencia de la cobertura de liquidez y distinguir el dato
observado de la proyección}. Es la única tarea del estudio que no terminó en éxito para todos: uno
de los dos participantes la resolvió en 76 segundos sin errores y el otro la resolvió de forma
parcial en 94 segundos con dos errores, porque leyó primero la proyección como si fuera el último
dato. La facilidad percibida cayó a 5 y 4 sobre 7, el valor más bajo de las seis tareas. Es el
hallazgo H1, el de mayor prioridad del estudio, y la recomendación R1 responde a él: rótulos
persistentes, fecha de corte y lenguaje inequívoco para la proyección.
\end{uxnota}

% ============================================================================
\clearpage
\subsection{Exportaciones}

\textbf{Para qué sirve.} Sacar del portal un conjunto de datos completo, en un archivo. Es la
rama 2.3 del módulo de exploración y tiene pantalla propia porque su ciclo de vida es largo: se
pide un trabajo, el trabajo corre en segundo plano y el resultado queda disponible durante una
ventana acotada. Pide perfil analista.

\textbf{Qué se ve.} Arriba, el formulario con \textbf{Conjunto de datos} ---Créditos, Liquidez o
Derivados---, \textbf{Formato} ---CSV u Hoja de Excel (XLSX)--- y \textbf{Filtros}. Abajo, el
\textbf{Historial de exportaciones}, que consulta el estado real de cada trabajo y se actualiza
solo mientras haya trabajos vivos. Una franja declara que el servidor estira ocho segundos la
duración de cada trabajo para que el estado intermedio pueda observarse, y que el trabajo, el
archivo y el enlace firmado sí son reales.

\figuraux[0.92\linewidth]{figuras/a4/tema/institucional_claro_6_exportacion.png}{Exportaciones con
sesión de perfil analista. El formulario declara el tope de la hoja de cálculo junto al control que
lo impone, la franja explica el retardo que la demostración añade a cada trabajo y el historial
compone cada tarjeta desde el estado real del trabajo, no desde un
guion.}{fig:a5-manual-exportacion}

\begin{uxpreg}
  \item Elige el conjunto de datos y el formato.
  \item Escribe los filtros si hacen falta, con la forma \texttt{columna=valor}: varios valores se
    separan con comas y varios filtros con punto y coma. Déjalo vacío para exportar el conjunto
    completo.
  \item Pulsa \textbf{Solicitar exportación}. La solicitud devuelve un identificador al instante y
    el trabajo pasa por los estados En cola, En proceso y Completado o Fallido.
  \item Sigue el trabajo en el historial y, cuando esté completo, usa \textbf{Descargar archivo}.
    El enlace está firmado y caduca a las 24 horas; pasado ese plazo se solicita la exportación de
    nuevo.
\end{uxpreg}

Un límite que conviene conocer antes de elegir formato: la hoja de Excel admite hasta 200\,000
filas y por encima de ese tope el trabajo termina fallido. El CSV no tiene tope y es el único
formato que saca el millón de filas del conjunto de Liquidez. La pantalla publica esa restricción
junto al propio control que la impone.

\begin{uxnota}[Ejemplo, tarea T3 de la prueba del 20 de agosto de 2026]
La tarea pedía \emph{solicitar la cartera de crédito en XLSX y localizar el estado del trabajo y la
vigencia del enlace}. Los dos participantes eligieron el conjunto Créditos y el formato Hoja de
Excel, enviaron la solicitud y siguieron el trabajo en el historial hasta ver su estado y su
caducidad. Los dos terminaron con éxito, en 64 y 79 segundos frente a una meta de 120, con un
error cada uno: revisaron las restricciones del formato después de haberlo elegido, no antes. Es
el hallazgo H3, y la recomendación R3 pide llevar tamaño, compatibilidad y uso recomendado al
propio selector.
\end{uxnota}

% ============================================================================
\clearpage
\subsection{Gobierno del dato}

\textbf{Para qué sirve.} Defender una cifra que ya se tiene. Exploración sirve para
\emph{encontrar} un dato y Gobierno para \emph{responder por él}: quién es su dueño, desde qué
vigencia, con qué certificación y por qué sistemas pasó antes de llegar a la pantalla.

\textbf{Qué se ve.} El \textbf{Diccionario de campos} con su buscador y su filtro por dominio.
Cada ficha trae el nombre físico, el nombre de negocio, la definición, la fuente, el responsable y
la vigencia, y abre el recorrido completo del dato. La pantalla no consulta nada hasta que se
escribe un término, así que abre sin pedirle campos al servidor. Al pie, el acceso cruzado a la
\textbf{Bitácora de accesos}, que declara que ese registro se administra desde el módulo de
Administración y por dónde se solicita.

\figuraux[0.92\linewidth]{figuras/a4/tema/institucional_claro_3_gobierno.png}{Gobierno del dato con
sesión de perfil analista, en su estado inicial. El diccionario explica que no consulta nada hasta
que se escribe un término, la barra lateral despliega las tres ramas del módulo ---diccionario y
metadatos, linaje y calidad, catálogo de fuentes--- y el acceso cruzado a la bitácora de accesos
nombra el módulo que la administra y la ruta por la que se solicita.}{fig:a5-manual-gobierno}

\begin{uxpreg}
  \item Escribe en \textbf{Campo cuyo origen hay que verificar} el nombre físico, el nombre de
    negocio o un sinónimo, con al menos dos caracteres.
  \item Acota por dominio si el término devuelve demasiadas fichas.
  \item Abre la ficha y lee el responsable, la vigencia y la certificación.
  \item Pulsa \textbf{Ver el linaje}. El panel dibuja el recorrido del dato por etapas ---origen,
    extracción, transformación, calidad y presentación--- y cada salto declara qué se le hizo al
    dato, quién responde por ese paso y entre qué fechas es válido.
\end{uxpreg}

\begin{uxnota}[Ejemplo, recorridos del 20 de agosto de 2026]
Gobierno del dato no tuvo tarea cronometrada propia en el estudio, de modo que aquí no hay tiempo
ni tasa de éxito que reportar. Sí hubo observación: los tres participantes de los recorridos
vinculados con su trabajo completaron sus actividades sin ninguna ayuda de la moderación, y los
tres reconocieron la relación entre campos y fuentes a partir de los metadatos. Es el hallazgo H6,
registrado como fortaleza en 3 de 3 personas expuestas, y la recomendación R5 pide conservar
visibles la fuente, el responsable y el estado por ese motivo. El recorrido que este manual
propone es la continuación natural de la tarea T2: una vez localizado el campo en Exploración, su
nombre físico abre el mismo panel de linaje que se consulta desde aquí.
\end{uxnota}

% ============================================================================
\clearpage
\subsection{Asistente conversacional}

\textbf{Para qué sirve.} Preguntar en lenguaje natural y seguir en vivo cada consulta que el
asistente ejecuta para responder. Es transversal a los cuatro módulos, así que se alcanza desde
cualquier pantalla, y pide perfil operativo o superior.

\textbf{Qué se ve.} Una franja que declara que las respuestas están guionizadas y que el
transporte, las tarjetas y la cancelación sí son reales. Debajo, la conversación; el campo
\textbf{Tu pregunta al asistente}; el estado de la generación; y el bloque \textbf{Contexto del
tablero}, que muestra lo que acompañaría a la pregunta si se hiciera en ese momento.

\figuraux[0.92\linewidth]{figuras/a4/tema/institucional_claro_7_asistente_resultado.png}{Un turno
resuelto del asistente. La tarjeta declara la consulta ejecutada y su duración, 0,3 segundos,
publica la cifra devuelta, muestra la tabla con la métrica y su valor, y nombra al pie el campo del
catálogo del que proviene. La respuesta escrita repite esa procedencia y declara que la cifra no es
una estimación del asistente.}{fig:a5-manual-asistente}

\begin{uxpreg}
  \item Escribe la pregunta y pulsa \textbf{Enviar pregunta}.
  \item Observa la \textbf{tarjeta de consulta}: el asistente anuncia la consulta antes de
    ejecutarla y publica su estado ---Anunciada, En ejecución, Resuelta o Fallida--- y su
    duración.
  \item Lee la cifra devuelta, la tabla de resultados y, al pie de la tarjeta, la
    \textbf{fuente del catálogo} de la que sale.
  \item Pulsa \textbf{Detener} si quieres abortar: la cancelación interrumpe la petición de
    verdad, no solo la vista.
  \item Abre \textbf{Detalle técnico} para ver qué herramienta se usó y sobre qué campo del
    catálogo.
\end{uxpreg}

El compromiso de esta pantalla es el más fuerte del producto y se puede comprobar a simple vista:
\textbf{ninguna cifra aparece sin una consulta detrás, y la consulta se enseña}. Una consulta que
falla no devuelve ninguna cifra, y la tarjeta lo dice con esas palabras.

\begin{uxnota}[Ejemplo, tarea T5 de la prueba del 20 de agosto de 2026]
La tarea pedía \emph{preguntar por liquidez e identificar qué consulta y qué fuente respaldan la
respuesta}. Los dos participantes escribieron la pregunta, esperaron a que la tarjeta pasara a
Resuelta y señalaron los dos elementos pedidos sin ninguna ayuda: la consulta anunciada en la
tarjeta y el campo del catálogo citado a su pie. Los dos terminaron con éxito, en 49 y 57 segundos
frente a una meta de 90, con cero errores y cero ayudas, y calificaron la facilidad con 6 sobre 7.
Es el hallazgo H5 y el estudio lo registra como fortaleza: la trazabilidad del asistente se
reconoce con facilidad.
\end{uxnota}

% ============================================================================
\clearpage
\subsection{Administración}

\textbf{Para qué sirve.} Administrar las cuentas del portal. Es el único módulo que pide perfil
administrador, y el portal comprueba ese rol en cada petición, no solo al pintar la pantalla.

\textbf{Qué se ve.} La tabla \textbf{Usuarios del portal}, con seis columnas: persona, correo,
rol, estado, última modificación y acciones. Sobre ella, un filtro por nombre, usuario o correo.
La barra lateral publica las cuatro ramas del módulo, y la propia pantalla declara que tres de
ellas ---solicitudes y aprobaciones, bitácora de accesos e integraciones--- llegan después de esta
entrega.

\figuraux[0.92\linewidth]{figuras/a4/tema/institucional_claro_4_administracion.png}{Administración
con sesión de perfil administrador. Cada fila publica la persona, su correo, su rol con el selector
que lo cambia, su estado, la fecha de su última modificación administrativa y la acción de
desactivar. La fila de la cuenta que abrió la sesión sustituye esa acción por la marca
\enquote{Es tu cuenta}.}{fig:a5-manual-administracion}

\begin{uxpreg}
  \item Escribe en el filtro para acotar la lista. El orden que la tabla recibe se conserva
    mientras nadie lo cambie, y cada columna anuncia por cuál está ordenada.
  \item Para cambiar el rol de una cuenta, elige el nuevo rol en el selector de su fila y confirma
    en el diálogo \enquote{¿Cambiar el rol de esta cuenta?}. El cambio se puede deshacer desde
    esta misma pantalla.
  \item Para retirar el acceso, pulsa \textbf{Desactivar} y confirma. La cuenta no se borra: queda
    desactivada y se reactiva desde aquí con el rol que ya tenía.
  \item Tu propia fila no ofrece esas dos acciones: en su lugar aparece la marca \textbf{Es tu
    cuenta}. El portal no permite que un administrador se degrade ni se desactive a sí mismo.
\end{uxpreg}

\begin{uxnota}[Ejemplo, tarea T6 de la prueba del 20 de agosto de 2026]
La tarea pedía \emph{cambiar al perfil administrador, localizar una cuenta y explicar cómo cambiar
o retirar su acceso sin modificarla}. Los dos participantes cambiaron de perfil desde la cabecera,
filtraron la lista, encontraron la cuenta y describieron las dos acciones sin ejecutar ninguna: la
meta de cero acciones accidentales se cumplió. Terminaron en 58 y 71 segundos frente a una meta de
90, con un error en uno de los dos casos ---buscó Administración antes de cambiar de perfil y tuvo
que volver a la cabecera---. Es el hallazgo H4 y la recomendación R4 pide confirmar el contexto al
cambiar de perfil, redirigiendo al inicio del nuevo perfil o avisando con un acceso directo a su
espacio.
\end{uxnota}

% ============================================================================
\clearpage
\subsection{Qué se ve cuando algo no sale}

Una pantalla que solo existe en su estado ideal no sirve para trabajar. El prototipo declara
cuatro estados no felices y los resuelve siempre de la misma manera, de modo que reconocer uno en
una pantalla basta para reconocerlo en todas.

\begin{uxtablalarga}{|L{2.6cm}|L{5.2cm}|Y|}{Los cuatro estados no felices: cómo se reconocen y qué hacer}{\thd{Estado} & \thd{Cómo se reconoce} & \thd{Qué hacer}}
  Cargando & Siluetas del mismo alto que el contenido que las sustituye: cinco filas de esqueleto en la tabla del catálogo, filas de esqueleto en la de usuarios, tarjetas de esqueleto en el historial de exportaciones & Esperar. La caja del esqueleto ocupa el mismo espacio que el contenido definitivo, así que nada se mueve de sitio al llegar el dato y ningún control cambia de posición bajo el puntero \\
  Vacío & Un título que dice que no hubo coincidencias y un cuerpo que explica por qué, no un espacio en blanco. La consulta llegó al servidor y respondió sin error & Seguir la sugerencia que da la propia pantalla: probar con el nombre de negocio en lugar del nombre físico, quitar el filtro de dominio o ampliar el rango de fechas \\
  Error & Un aviso que conserva lo que ya estaba escrito ---el término de búsqueda, los filtros aplicados--- y ofrece \textbf{Reintentar} & Reintentar sin volver a teclear. En Inicio el fallo queda contenido en el bloque que lo sufre: un fallo del buscador no vacía las alertas \\
  Sin permiso & Un aviso que dice qué perfil pide la pantalla y cuál es el tuyo, dibujado donde habría estado la pantalla y sin cambiar la dirección. No hay botón de reintentar & Solicitar el acceso en Administración, rama de solicitudes y aprobaciones, o volver a tu espacio de trabajo con el único enlace que el aviso ofrece \\
\end{uxtablalarga}

Dos precisiones sobre el último estado, porque son deliberadas. La primera: \textbf{no hay control
de reintento}, y la ausencia es intencionada ---reintentar un rechazo no cambia el rol de nadie, y
ofrecer el botón enseñaría a insistir contra una puerta que no va a abrirse---. La segunda: el
aviso no siempre ocupa la pantalla entera. La ruta de exploración no exige permiso, así que ahí el
rechazo recae sobre sus dos continuaciones gobernadas, que se nombran, declaran qué perfil piden y
quedan visibles sin ser accionables.

La documentación de la Actividad 4 publica estos cuatro estados pantalla por pantalla, en una
matriz de siete filas por cuatro columnas sin ninguna celda en blanco, y distingue los que están
construidos de los que están especificados con su regla de comportamiento. Este manual describe lo
que se ve; aquella matriz declara dónde se ve.

% ============================================================================
\clearpage
\subsection{Preguntas frecuentes}

\begin{uxlist}
  \lead{¿Los datos son reales?}{No. Todo el contenido sale de generadores sintéticos con semilla
    fija. El portal lo declara en la franja de alcance de las diez rutas, en el distintivo de cada
    bloque con muestras y en el bloque de procedencia del tablero, que publica el archivo de
    origen y la semilla.}
  \lead{¿Por qué la primera pantalla tarda más que las siguientes?}{Porque el servicio se apaga
    cuando nadie lo usa y arranca con la primera petición. Ese arranque en frío se midió el 22 de
    agosto de 2026 en 2.4 segundos; las peticiones siguientes ya no lo pagan. Es la contrapartida
    de un despliegue que no cobra por estar encendido sin tráfico.}
  \lead{¿Por qué el asistente siempre muestra la consulta y la fuente?}{Porque ninguna cifra puede
    aparecer sin una consulta detrás. La tarjeta anuncia la consulta antes de ejecutarla, publica
    su estado y su duración, y cita al pie el campo del catálogo del que sale el número; la
    respuesta escrita repite esa procedencia. Una consulta que falla no devuelve ninguna cifra, y
    la tarjeta lo dice. Es lo contrario de pedirle confianza al lector: se le entrega el rastro
    para que la compruebe.}
  \lead{¿Cómo distingo un dato de una proyección?}{Por tres marcas simultáneas: la línea
    proyectada es discontinua y su marcador hueco, la leyenda nombra los dos tramos y un rótulo
    declara cuál fue el último mes observado. Además, cada tarjeta publica su método y su
    R\textsuperscript{2}, y advierte que detrás no hay aprendizaje automático sino una recta
    ajustada a la historia sintética.}
  \lead{¿Puedo cambiar de perfil sin cerrar sesión?}{Sí. El selector de perfil de la cabecera abre
    los cuatro y cambia la sesión en el sitio. Con el perfil cambia lo que cada pantalla muestra,
    empezando por la composición de Inicio.}
  \lead{¿Cuánto dura la sesión?}{Treinta minutos. Al caducar, el portal devuelve a la pantalla de
    acceso, explica lo ocurrido y despliega solo el formulario de credenciales.}
  \lead{¿Por qué no veo el módulo de Administración?}{Porque ese módulo pide perfil administrador.
    Entra con el perfil de demostración \textbf{Administración} desde la pantalla de acceso o
    desde el selector de la cabecera.}
  \lead{¿Por qué mi exportación en Excel terminó fallida?}{Porque la hoja de cálculo admite hasta
    200\,000 filas y el conjunto pedido las supera. Acota con filtros o pide la exportación en
    CSV, que no tiene tope.}
  \lead{¿Cuánto vive el enlace de descarga?}{Veinticuatro horas. Pasado ese plazo el enlace ya no
    sirve y la exportación se solicita de nuevo, lo que produce un enlace nuevo.}
  \lead{¿El portal está en español o en inglés?}{En los dos. Los botones \textbf{ES} y \textbf{EN}
    de la cabecera cambian el idioma de toda la interfaz sin perder la pantalla en la que estás,
    porque la dirección es la misma en ambos idiomas.}
  \lead{¿Se puede usar solo con el teclado?}{Sí. La primera parada del tabulador en cada pantalla
    es el enlace \enquote{Saltar al contenido}, las columnas ordenables anuncian su estado a los
    lectores de pantalla y toda gráfica tiene una alternativa en tabla y un resumen escrito
    equivalente.}
\end{uxlist}
